\section*{Kapitel 4 --- CAN, registerkartan och arkitekturen}
\addcontentsline{toc}{section}{Kapitel 4 --- CAN, registerkartan och arkitekturen}

\solution{ex:4:1}{Vem gör vad}
Ramsynkronisering, längdfält, checksumma, kvittens och kollisionshantering görs alla av
\textbf{hårdvaran}: SOF och EOF, DLC, CRC-15, ACK-biten inuti framen, och arbitreringen som bussen
löser elektriskt.

De två raderna utan självklart svar är \textbf{adressering} och \textbf{sekvensnummer}. CAN har
ingen destinationsadress och inget sekvensnummer alls: identifieraren beskriver innehållet, och
behövs ett löpnummer får applikationen lägga det i nyttolasten.

\textbf{Timeout och omsändning} samt \textbf{dubblettdetektering} är ett tredje fall: en
fullständig CAN-kontroller sköter omsändningen i hårdvara, men den förenklade i den här
konstruktionen gör det inte --- se \secref{sec:limits}.

\solution{ex:4:2}{Packa en frame}
Databyte 0 ligger i de mest signifikanta bitarna, och \code{TX_DATA_LO} bär de \emph{första} fyra
bytesen:

\begin{console}
TX_ID      = 0x000002A7
TX_DLC     = 0x00000005
TX_DATA_LO = 0x01020304
TX_DATA_HI = 0x05000000
TX_SEND    = 0x00000001
\end{console}

Ordningen spelar roll därför att skrivningen till \code{TX_SEND} inte lagrar ett värde utan
\emph{utlöser en sändning}: kontrollern läser då de övriga registren. Skrivs triggern först skickas
föregående frames innehåll.

\solution{ex:4:3}{Packa upp en frame}
\ans{a} \code{0x481}.
\ans{b} \code{DE AD BE} --- tre bytes, eftersom \code{RX_DLC} är 3 och \code{RX_DATA_LO} bär byte
0--3 med byte 0 överst.
\ans{c} Noll. Drivrutinen ska maskera mot \code{dlc}; bytes ovanför den mottagna längden är inte
framens innehåll, även om hårdvaran råkar nollställa dem.
\ans{d} \code{RX_ACK}, med \code{0x1}. Bit 0 måste vara satt.

\solution{ex:4:4}{Fyra påståenden}
\textbf{a, b och c beskriver buggar.}

\ans{a} Bit 0 betyder \emph{porten är fri igen}, inte \emph{framen kom fram} --- en förlorad
arbitrering ger tillbaka samma bit. Läs \code{ERROR_FLAGS} efteråt.
\ans{b} \code{ERROR_FLAGS} nollställs bara av ett \emph{helt nollställt} ord. \code{0x1} gör
ingenting, så felbiten går aldrig ned.
\ans{c} Bytes ovanför \code{RX_DLC} hör inte till framen. Att kopiera alla åtta rapporterar upp
till sex bytes som aldrig sändes.
\ans{d} Korrekt. Utan kontrollen kan en skrivning ske mitt i en pågående sändning.

\solution{ex:4:5}{Lagergränsen}
\ans{a} Falskt. \code{EchoNode} ligger ovanför interfacet och får inte känna till någon konkret
drivrutin.
\ans{b} Sant. \code{Spi} \emph{är} en \code{Interface} och måste inkludera den.
\ans{c} Falskt. Stubben simulerar CAN i minnet och har ingen transport.
\ans{d} Falskt. Transporten flyttar bytes och vet ingenting om register --- det är hela poängen
med att den är dum.
\ans{e} Falskt. \code{EchoNode} testas genom \code{Stub}, som inte har någon transport.
